% Options for packages loaded elsewhere
% Options for packages loaded elsewhere
\PassOptionsToPackage{unicode}{hyperref}
\PassOptionsToPackage{hyphens}{url}
%
\documentclass[
  ignorenonframetext,
  aspectratio=169,
  russian,
]{beamer}
\newif\ifbibliography
\usepackage{pgfpages}
\setbeamertemplate{caption}[numbered]
\setbeamertemplate{caption label separator}{: }
\setbeamercolor{caption name}{fg=normal text.fg}
\beamertemplatenavigationsymbolshorizontal
% remove section numbering
\setbeamertemplate{part page}{
  \centering
  \begin{beamercolorbox}[sep=16pt,center]{part title}
    \usebeamerfont{part title}\insertpart\par
  \end{beamercolorbox}
}
\setbeamertemplate{section page}{
  \centering
  \begin{beamercolorbox}[sep=12pt,center]{section title}
    \usebeamerfont{section title}\insertsection\par
  \end{beamercolorbox}
}
\setbeamertemplate{subsection page}{
  \centering
  \begin{beamercolorbox}[sep=8pt,center]{subsection title}
    \usebeamerfont{subsection title}\insertsubsection\par
  \end{beamercolorbox}
}
% Prevent slide breaks in the middle of a paragraph
\widowpenalties 1 10000
\raggedbottom
\AtBeginPart{
  \frame{\partpage}
}
\AtBeginSection{
  \ifbibliography
  \else
    \frame{\sectionpage}
  \fi
}
\AtBeginSubsection{
  \frame{\subsectionpage}
}
\usepackage{iftex}
\ifPDFTeX
  \usepackage[T1]{fontenc}
  \usepackage[utf8]{inputenc}
  \usepackage{textcomp} % provide euro and other symbols
\else % if luatex or xetex
  \usepackage{unicode-math} % this also loads fontspec
  \defaultfontfeatures{Scale=MatchLowercase}
  \defaultfontfeatures[\rmfamily]{Ligatures=TeX,Scale=1}
\fi
\usepackage{lmodern}

\ifPDFTeX\else
  % xetex/luatex font selection
\fi
% Use upquote if available, for straight quotes in verbatim environments
\IfFileExists{upquote.sty}{\usepackage{upquote}}{}
\IfFileExists{microtype.sty}{% use microtype if available
  \usepackage[]{microtype}
  \UseMicrotypeSet[protrusion]{basicmath} % disable protrusion for tt fonts
}{}


\usepackage{longtable,booktabs,array}
\usepackage{calc} % for calculating minipage widths
\usepackage{caption}
% Make caption package work with longtable
\makeatletter
\def\fnum@table{\tablename~\thetable}
\makeatother
\usepackage{graphicx}
\makeatletter
\newsavebox\pandoc@box
\newcommand*\pandocbounded[1]{% scales image to fit in text height/width
  \sbox\pandoc@box{#1}%
  \Gscale@div\@tempa{\textheight}{\dimexpr\ht\pandoc@box+\dp\pandoc@box\relax}%
  \Gscale@div\@tempb{\linewidth}{\wd\pandoc@box}%
  \ifdim\@tempb\p@<\@tempa\p@\let\@tempa\@tempb\fi% select the smaller of both
  \ifdim\@tempa\p@<\p@\scalebox{\@tempa}{\usebox\pandoc@box}%
  \else\usebox{\pandoc@box}%
  \fi%
}
% Set default figure placement to htbp
\def\fps@figure{htbp}
\makeatother



\ifLuaTeX
\usepackage[bidi=basic,provide=*]{babel}
\else
\usepackage[bidi=default,provide=*]{babel}
\fi
% get rid of language-specific shorthands (see #6817):
\let\LanguageShortHands\languageshorthands
\def\languageshorthands#1{}


\setlength{\emergencystretch}{3em} % prevent overfull lines

\providecommand{\tightlist}{%
  \setlength{\itemsep}{0pt}\setlength{\parskip}{0pt}}



 

\usepackage[]{csquotes}

\IfFileExists{plex-otf.sty}{
  %% Full TeXlive
  % \usepackage[%
  %   % math,
  %   RM={Scale=0.94},SS={Scale=0.94},SScon={Scale=0.94},TT={Scale=MatchLowercase,FakeStretch=0.9},DefaultFeatures={Ligatures=Common}
  % ]{plex-otf}
}{
  %% TinyTeX
  \usepackage{libertine}
}

%%% Load theme
% https://deic.uab.cat/~iblanes/beamer_gallery/
\IfFileExists{beamerthemegotham.sty}{
  %% Full TeXlive
  \usetheme{gotham}
  \gothamset{
    numbering=totalpagenumber,
    parttocframe default=off,
    sectiontocframe default=off,
    subsectiontocframe default=off,
  }
}{
  %% TinyTeX
  \usetheme{Madrid}
}
\makeatletter
\@ifpackageloaded{caption}{}{\usepackage{caption}}
\AtBeginDocument{%
\ifdefined\contentsname
  \renewcommand*\contentsname{Содержание}
\else
  \newcommand\contentsname{Содержание}
\fi
\ifdefined\listfigurename
  \renewcommand*\listfigurename{Список иллюстраций}
\else
  \newcommand\listfigurename{Список иллюстраций}
\fi
\ifdefined\listtablename
  \renewcommand*\listtablename{Список таблиц}
\else
  \newcommand\listtablename{Список таблиц}
\fi
\ifdefined\figurename
  \renewcommand*\figurename{Рисунок}
\else
  \newcommand\figurename{Рисунок}
\fi
\ifdefined\tablename
  \renewcommand*\tablename{Таблица}
\else
  \newcommand\tablename{Таблица}
\fi
}
\@ifpackageloaded{float}{}{\usepackage{float}}
\floatstyle{ruled}
\@ifundefined{c@chapter}{\newfloat{codelisting}{h}{lop}}{\newfloat{codelisting}{h}{lop}[chapter]}
\floatname{codelisting}{Список}
\newcommand*\listoflistings{\listof{codelisting}{Листинги}}
\makeatother
\makeatletter
\makeatother
\makeatletter
\@ifpackageloaded{caption}{}{\usepackage{caption}}
\@ifpackageloaded{subcaption}{}{\usepackage{subcaption}}
\makeatother

\usepackage{bookmark}
\IfFileExists{xurl.sty}{\usepackage{xurl}}{} % add URL line breaks if available
\urlstyle{same}
\hypersetup{
  pdftitle={Лабораторная работа №13},
  pdfauthor={Ко Антон Геннадьевич},
  pdflang={ru-RU},
  hidelinks,
  pdfcreator={LaTeX via pandoc}}


\title{Лабораторная работа №13}
\subtitle{Статическая маршрутизация в Интернете. Планирование}
\author{Ко Антон Геннадьевич}
\date{2026-05-09}

\begin{document}
\frame{\titlepage}

\renewcommand*\contentsname{Содержание}
\begin{frame}[allowframebreaks]
  \frametitle{Содержание}
  \setcounter{tocdepth}{1}
  \tableofcontents
\end{frame}
\setcounter{tocdepth}{1}
\tableofcontents
}

\section{Информация}\label{ux438ux43dux444ux43eux440ux43cux430ux446ux438ux44f}

\begin{frame}{Докладчик}
\phantomsection\label{ux434ux43eux43aux43bux430ux434ux447ux438ux43a}
\begin{columns}[c]
\begin{column}{0.7\linewidth}
\begin{itemize}[<+->]
\tightlist
\item
  Ко Антон Геннадьевич
\item
  студент
\item
  Российский университет дружбы народов им. П. Лумумбы
\item
  \href{mailto:1132221551@rudn.ru}{\nolinkurl{1132221551@rudn.ru}}
\item
  \url{https://SenDerMen04.github.io/ru/}
\end{itemize}
\end{column}

\begin{column}{0.3\linewidth}
\pandocbounded{\includegraphics[keepaspectratio]{./image/me.jpg}}
\end{column}
\end{columns}
\end{frame}

\begin{frame}{Цель работы}
\phantomsection\label{ux446ux435ux43bux44c-ux440ux430ux431ux43eux442ux44b}
Провести подготовительные мероприятия по организации взаимодействия
через сеть провайдера посредством статической маршрутизации локальной
сети с сетью основного здания, расположенного в 42-м квартале в Москве,
и сетью филиала, расположенного в г. Сочи.
\end{frame}

\begin{frame}{Выполнение работы}
\phantomsection\label{ux432ux44bux43fux43eux43bux43dux435ux43dux438ux435-ux440ux430ux431ux43eux442ux44b}
На схеме нашего проекта разместим необходимое оборудование: 4
медиаконвертера (Repeater-PT), 2 маршрутизатора типа Cisco 2811, 1
маршрутизирующий коммутатор типа Cisco 3560-24PS, 2 коммутатора типа
Cisco 2950-24, коммутатор Cisco 2950-24T, 3 оконечных устройства типа
PC-PT. А также присвоим им названия и проведём соединение объектов
согласно скорректированной нами схеме (рис. \#fig:007):

\begin{figure}

{\centering \includegraphics[width=1\linewidth,height=\textheight,keepaspectratio]{image/7.png}

}

\caption{Размещение необходимого оборудования, присвоение названий и
соединение объектов}

\end{figure}%
\end{frame}

\begin{frame}[fragile]
На медиаконвертерах заменим имеющиеся модули на
\texttt{PT-REPEATERNM-1FFE} и \texttt{PT-REPEATER-NM-1CFE} для
подключения витой пары по технологии Fast Ethernet и оптоволокна
соответственно (рис. \#fig:008):

\begin{figure}

{\centering \includegraphics[width=1\linewidth,height=\textheight,keepaspectratio]{image/8.png}

}

\caption{Замена модулей на медиаконвертерах}

\end{figure}%
\end{frame}

\begin{frame}[fragile]
Далее на маршрутизаторе \texttt{msk-q42-agko-gw-1} добавим
дополнительный интерфейс NM-2FE2W (рис. \#fig:009):

\begin{figure}

{\centering \includegraphics[width=1\linewidth,height=\textheight,keepaspectratio]{image/9.png}

}

\caption{Добавление на маршрутизаторе msk-q42-agko-gw-1 дополнительного
интерфейса NM-2FE2W}

\end{figure}%
\end{frame}

\begin{frame}
В физической рабочей области Packet Tracer добавим в г. Москва здание
42-го квартала и присвоим ему название (рис. \#fig:010):

\begin{figure}

{\centering \includegraphics[width=1\linewidth,height=\textheight,keepaspectratio]{image/10.png}

}

\caption{Добавление в физической рабочей области в г. Москва здания
42-го квартала}

\end{figure}%
\end{frame}

\begin{frame}
Затем в физической рабочей области добавим город Сочи и в нём здание
филиала, присвоим ему соответствующее название (рис. \#fig:011):

\begin{figure}

{\centering \includegraphics[width=1\linewidth,height=\textheight,keepaspectratio]{image/11.png}

}

\caption{Добавление города Сочи и здания филиала}

\end{figure}%
\end{frame}

\begin{frame}
После чего нужно перенести из сети «Донская» оборудование сети 42-го
квартала и сети филиала в соответствующие здания и выполнить расстановку
(рис. \#fig:012 -- \#fig:014):

\begin{figure}

{\centering \includegraphics[width=1\linewidth,height=\textheight,keepaspectratio]{image/12.png}

}

\caption{Перенос оборудования сети 42-го квартала и сети филиала в
соответствующие здания}

\end{figure}%
\end{frame}

\begin{frame}
\begin{figure}

{\centering \includegraphics[width=1\linewidth,height=\textheight,keepaspectratio]{image/13.png}

}

\caption{Размещение объектов в основном здании 42-го квартала в Москве}

\end{figure}%
\end{frame}

\begin{frame}
\begin{figure}

{\centering \includegraphics[width=1\linewidth,height=\textheight,keepaspectratio]{image/14.png}

}

\caption{Размещение объектов в здании филиала в г. Сочи}

\end{figure}%
\end{frame}

\begin{frame}
На последнем шаге выполним первоначальную настройку оборудования (рис.
\#fig:015 -- \#fig:020):

\begin{figure}

{\centering \includegraphics[width=1\linewidth,height=\textheight,keepaspectratio]{image/15.png}

}

\caption{Первоначальная настройка маршрутизатора msk-q42-agko-gw-1}

\end{figure}%
\end{frame}

\begin{frame}
\begin{figure}

{\centering \includegraphics[width=1\linewidth,height=\textheight,keepaspectratio]{image/16.png}

}

\caption{Первоначальная настройка коммутатора msk-q42-agko-sw-1}

\end{figure}%
\end{frame}

\begin{frame}
\begin{figure}

{\centering \includegraphics[width=1\linewidth,height=\textheight,keepaspectratio]{image/17.png}

}

\caption{Первоначальная настройка маршрутизирующего коммутатора
msk-hostel-agko-gw-1}

\end{figure}%
\end{frame}

\begin{frame}
\begin{figure}

{\centering \includegraphics[width=1\linewidth,height=\textheight,keepaspectratio]{image/18.png}

}

\caption{Первоначальная настройка коммутатора msk-hostel-sw-1}

\end{figure}%
\end{frame}

\begin{frame}
\begin{figure}

{\centering \includegraphics[width=1\linewidth,height=\textheight,keepaspectratio]{image/19.png}

}

\caption{Первоначальная настройка коммутатора sch-sochi-agko-sw-1}

\end{figure}%
\end{frame}

\begin{frame}
\begin{figure}

{\centering \includegraphics[width=1\linewidth,height=\textheight,keepaspectratio]{image/20.png}

}

\caption{Первоначальная настройка маршрутизатора sch-sochi-agko-gw-1}

\end{figure}%
\end{frame}

\begin{frame}{Вывод}
\phantomsection\label{ux432ux44bux432ux43eux434}
В ходе выполнения лабораторной работы мы провели подготовительные
мероприятия по организации взаимодействия через сеть провайдера
посредством статической маршрутизации локальной сети с сетью основного
здания, расположенного в 42-м квартале в Москве, и сетью филиала,
расположенного в г. Сочи.
\end{frame}




\end{document}
